\documentclass[11pt]{article}

% Change "review" to "final" to generate the final (sometimes called camera-ready) version.
% Change to "preprint" to generate a non-anonymous version with page numbers.
\usepackage[review]{acl}

% Standard package includes
\usepackage{times}
\usepackage{latexsym}

% For proper rendering and hyphenation of words containing Latin characters (including in bib files)
\usepackage[T1]{fontenc}
% For Vietnamese characters
% \usepackage[T5]{fontenc}
% See https://www.latex-project.org/help/documentation/encguide.pdf for other character sets

% This assumes your files are encoded as UTF8
\usepackage[utf8]{inputenc}

% This is not strictly necessary, and may be commented out,
% but it will improve the layout of the manuscript,
% and will typically save some space.
\usepackage{microtype}

% This is also not strictly necessary, and may be commented out.
% However, it will improve the aesthetics of text in
% the typewriter font.
% Commented out: inconsolata.sty is not available in this environment's TeX install.
% \usepackage{inconsolata}

%Including images in your LaTeX document requires adding
%additional package(s)
\usepackage{graphicx}

% If the title and author information does not fit in the area allocated, uncomment the following
%
%\setlength\titlebox{<dim>}
%
% and set <dim> to something 5cm or larger.

% \usepackage[margin=1in]{geometry}
\usepackage{booktabs}
\usepackage{multirow}
\usepackage{amsmath}
\usepackage{natbib}
\usepackage{xcolor}

% Skeleton conventions:
%  - Every section carries \itemize bullets = the meat to be expanded into prose.
%  - \todo{...} = decision or number still needed before writing the prose.
%  - Single-table constraint: tab:main-results is the ONLY table in the paper.
%    All ablation/analysis numbers go into prose deltas or figures.
\newcommand{\todo}[1]{\textcolor{red}{[TODO: #1]}}
\newcommand{\method}{SCOTECI}

\title{\method{}: Synthetic Chain-of-Thought Demonstrations for Document-Level Event Causality Identification\\
{\normalsize \todo{working title --- alternatives: ``Documents, not Pairs: ...'' / ``Rationalized Few-Shot Prompting for Document-Level ECI''}}}
\author{}
\date{}

\begin{document}

\maketitle

% ======================================================================
\begin{abstract}
Large language models (LLMs) applied to Event Causality Identification (ECI) exhibit
\emph{causal hallucination}: near-total recall paired with unusable precision
\citep{gao2023chatgpt}. Prior LLM-based scaffolds for ECI, such as Dr.ECI and SERE,
fight this at the \emph{pair} level, issuing one call per candidate event pair with
bare input-to-label demonstrations \citep{cai2025dreci, hao2026sere}. We propose
\method{}, which instead (i) reformulates LLM-based ECI at the \emph{document} level,
labelling every pair of a document jointly in a single call under a structured
reasoning protocol with explicit coherence rules, and (ii) equips its few-shot
demonstrations with \emph{synthetic chain-of-thought}: gold-guided reasoning generated
by an LLM and then decontaminated so it reads as genuine discovery rather than answer
confirmation. The method needs no fine-tuning, no human-written rationales, and no
external knowledge base. On Causal-TimeBank, \method{} reaches F1 48.1 (precision
61.1) against F1 20.0 for the best prior LLM method --- a +28.1 F1 gain obtained by
breaking out of the degenerate high-recall/low-precision regime rather than trading
recall for precision, on the dataset where positives are sparsest and hallucination is
costliest. It also attains the best precision on MAVEN-ERE (40.8) and matches SERE's
F1 on EventStoryLine (51.7 vs.\ 49.9). Beyond the headline numbers, we present an
evidence-driven analysis of the method itself: majority-vote aggregation over
demonstration ensembles is empirically F1-optimal, precision \emph{rises} rather than
falls with document causal density, and synthetic chain-of-thought demonstrations help
chat-tuned backbones but measurably hurt a reasoning-tuned one.
\end{abstract}

% ======================================================================
\section{Introduction}
\label{sec:intro}

Event Causality Identification (ECI) asks, given a document and its marked event
mentions, whether a directed causal relation holds between each pair of events
\citep{cheng2025survey}. It is a building block for event forecasting, causal
question answering, and narrative and discourse understanding, and has been pursued
along two largely disjoint paths: fine-tuned encoders and, more recently, prompted
LLMs. Fine-tuned encoders --- BERT/RoBERTa-era models augmented with event-relational
graphs or external knowledge --- are strong in-domain but require per-corpus
annotation and generalize poorly outside the corpus they were trained on
\citep{cai2025dreci, hao2026sere}. Prompted LLMs sidestep the annotation cost
entirely, but exhibit \emph{causal hallucination}: on Causal-TimeBank, a plain
zero-shot or chain-of-thought prompt reaches 80--96\% recall with precision below
10\% (Table~\ref{tab:main-results}, Base/CoT rows) --- the model asserts a causal
link almost anywhere a temporal or narrative connection exists \citep{gao2023chatgpt}.

The LLM scaffolds designed to fix this --- Dr.ECI's four-agent discovery/reasoning
pipeline \citep{cai2025dreci}, SERE's structurally retrieved demonstrations over
ConceptNet paths and syntax-tree edit distance \citep{hao2026sere}, and DICP's
learned deep in-context prompts \citep{mu2025dicp} --- all operate at the
\emph{pair} level: one LLM call per candidate event pair, re-paying the full
document context every time. This has three consequences. First, cost scales as
$O(|\text{pairs}|)$ calls per document ($O(n^2)$ in the number of mentions), each
of which resends the whole document. Second, each pair is judged in isolation, so
nothing enforces symmetry or transitivity across the pairs of the same document, and
the model never sees the document's causal structure as a whole. Third, these
scaffolds were engineered around the specific weaknesses of 2023--2024-era models;
as base models improve, the scaffold increasingly caps rather than unlocks their
capability.

We propose \method{}, built on two ideas that invert this design. The first is to
operate at the \emph{document} level: a single call labels every pair of a document
jointly, guided by an in-prompt reasoning protocol (mention pre-filter $\to$
per-pair arguments for and against $\to$ plausibility ranking $\to$ final JSON)
together with explicit coherence rules --- symmetry and dataset-specific transitivity
--- that are checkable only because the whole graph is produced at once. This shifts
the burden of consistency from external orchestration to the model's own long-context
reasoning, which is a deliberate bet: the method should get better, not merely
cheaper, as the underlying model improves.

The second idea addresses what document-level demonstrations should look like. Bare
input$\to$gold-JSON demonstrations teach the model the output format but not how to
reason about it, and document-level human-written rationales for ECI do not exist
and would be expensive to produce. Since gold labels for training documents
\emph{are} available, we generate the reasoning instead: a demonstration document is
shown to an LLM together with its own task prompt and its gold labels, and the model
is asked to write reasoning that leads naturally to those labels (\emph{gold-guided
generation}); the response is then rewritten from scratch to remove any phrasing that
betrays foreknowledge of the answer (\emph{decontamination}), so that every sentence
reads as discovery rather than confirmation. An optional three-step variant inserts a
blind attempt before the gold-guided pass, preserving more of the texture of a
genuine, imperfect first try. Each demonstration becomes a (user, assistant) exchange
in the prompt, so the model effectively ``remembers having solved'' a handful of
documents correctly, reasoning included.

On Causal-TimeBank, \method{} reaches F1 48.1 (precision 61.1, recall 39.6) against
F1 20.0 for the best prior LLM method (SERE) --- a +28.1 F1 gain, with precision more
than four times any LLM baseline in the comparison (Table~\ref{tab:main-results}).
This is the dataset where positives are sparsest and hallucination is most punishing,
and it is also the dataset where the balanced precision/recall regime most clearly
replaces the degenerate high-recall/near-zero-precision one that every baseline
exhibits. \method{} additionally attains the best precision on MAVEN-ERE (40.8) and
an F1 that edges out SERE on EventStoryLine (51.7 vs.\ 49.9). All three \method{}
numbers are reported on the intra-sentence subset of pairs, which is the only subset
SERE itself evaluates, so this is the fair like-for-like comparison; full-document
(intra- and inter-sentence) prediction is a capability the document-level
formulation affords --- one call covers all pairs regardless of how many of them
cross a sentence boundary, whereas a pair-level method pays for every additional
inter-sentence pair as another call --- but is not yet measured for \method{} and is
left to future work (see Limitations).

Beyond the headline table, we use \method{}'s logged predictions to answer questions
the original design left open. We show that the majority-vote aggregation used
across resampled passes is empirically F1-optimal rather than merely
``precision-preserving'' as assumed, that precision \emph{rises} rather than falls
with a document's causal density --- the opposite of the attention-dilution failure
mode we originally expected --- and that synthetic chain-of-thought demonstrations
help chat-tuned backbones but measurably hurt a reasoning-tuned one, a finding with
direct implications for which future backbones the method should target.

Our contributions are:
\begin{itemize}
  \item A document-level reformulation of LLM-based ECI, with a structured in-prompt
        reasoning protocol and explicit coherence rules, that issues $O(1)$ LLM
        calls per document instead of $O(|\text{pairs}|)$ and produces predictions
        that are consistent across a document's pairs by construction.
  \item A synthetic chain-of-thought demonstration pipeline --- gold-guided
        generation followed by decontamination --- that is model-agnostic and needs
        no human-written rationales or external knowledge base.
  \item State-of-the-art LLM results on Causal-TimeBank by a wide margin, with
        balanced precision/recall on all three evaluated datasets, plus an
        evidence-based analysis of aggregation, document density, backbone
        interaction, cost, and robustness that goes beyond the headline comparison.
\end{itemize}

% ======================================================================
\section{Related Work}
\label{sec:related}


\subsection{Supervised and Knowledge-Enhanced ECI}
Before LLM prompting, ECI was dominated by fine-tuned encoders that inject
structural or external signal into a BERT/RoBERTa-style backbone: event-relational
graph transformers that propagate information across a document's mention graph
\citep{chen2022ergo}, sparse event representations for discriminative document-level
event-relation extraction \citep{yuan2023sendir}, joint identify-while-learn
objectives that couple mention and relation supervision \citep{liu2024ilif},
knowledge-guided binary question-answering formulations of the pairwise decision
\citep{wang2024knowqa}, and heterogeneous graph contrastive transfer for zero-shot
cross-lingual ECI \citep{he2024gimc}. These remain the strongest in-domain systems on
their respective benchmarks, but each requires gold-annotated training data for the
specific corpus (often the specific label scheme) it is evaluated on, and transfer
across corpora or languages is itself an active research problem rather than a given
\citep{he2024gimc}. \method{} requires no task-specific training at all: the
training split is used only as a pool of few-shot demonstrations.

\subsection{LLM-based ECI}
\citet{gao2023chatgpt} first showed systematically that ChatGPT-class models
over-predict causal links --- causal hallucination --- and established the
preprocessing and evaluation protocol that we and our baselines follow. Two
scaffolds subsequently targeted this failure mode at the pair level. Dr.ECI
decomposes each candidate pair into a discovery stage (a Causal Explorer and a
Mediator Detector) and a reasoning stage (Direct and Indirect Reasoners), guided by
causal-pattern DAG priors \citep{cai2025dreci}. SERE retrieves structurally similar
pair-level examples using ConceptNet shortest-path edit distance, dependency-tree
edit distance, and causal-pattern filtering, and balances positive and negative
demonstrations by construction \citep{hao2026sere}. Closer to our synthetic-rationale
idea, \citet{wang2024synthetic} frame ECI through a synthetic-control,
counterfactual-estimation lens, and DICP learns deep in-context prompts per pair
rather than retrieving them \citep{mu2025dicp}. All of these methods (i) prompt once
per pair and (ii) show label-only demonstrations. \method{} prompts once per
document and shows demonstrations that carry explicit reasoning.

\subsection{In-Context Learning and Synthetic Rationales}
Chain-of-thought prompting improves multi-step reasoning by asking a model to
externalize intermediate steps before its final answer \citep{wei2022chain},
building on the broader observation that large models are effective few-shot
learners given the right demonstrations \citep{brown2020language}. In-context
learning, however, is highly sensitive to which demonstrations are shown and how
they are phrased \citep{liu2022makes, min2022rethinking}. Where rationales are not
naturally available, prior work has synthesized them: Auto-CoT clusters queries and
applies zero-shot CoT to generate rationales for demonstration selection
\citep{zhang2023autocot}; STaR bootstraps rationales through
\emph{rationalization} --- generating reasoning conditioned on a known answer ---
but consumes the result as fine-tuning data \citep{zelikman2022star}; step-by-step
distillation similarly extracts rationales as a training signal for smaller student
models \citep{hsieh2023distilling}. \method{} uses the same rationalization trick as
STaR, but differently: the synthesized reasoning is consumed \emph{at inference
time} as in-context demonstrations rather than as fine-tuning data, it operates at
document rather than sentence/example scale, and it adds an explicit decontamination
pass to strip answer-leakage phrasing --- without which, we find, demonstrations
teach the model to assert conclusions rather than to reason toward them
(\S\ref{sec:ablations}).

% ======================================================================
\section{Task Formulation}
\label{sec:task}

Given a document $D$ with event mentions marked inline in \texttt{<ID word>}
format, $M = \{m_1, \dots, m_n\}$, and a candidate pair set $P$ --- the annotated
pair list for corpora that provide one (the \emph{constrained} setting, used for
MECI and Causal-TimeBank) or all unordered mention pairs otherwise (the
\emph{open-ended} setting, used for MAVEN-ERE and EventStoryLine) --- the task is to
produce a complete labelling $\ell: P \to \{\textit{causes}, \textit{causedby},
\textit{norel}\}$: a unified, directed three-way label space into which every
corpus' native annotation scheme is mapped before evaluation. Concretely: MECI's
\textsc{CauseEffect}/\textsc{EffectCause}/\textsc{NoRel} labels map one-to-one onto
\textit{causes}/\textit{causedby}/\textit{norel}; MAVEN-ERE's \textsc{CAUSE} and
\textsc{PRECONDITION} relation types are both collapsed onto the same directed
causes/causedby pair; EventStoryLine's \textsc{PRECONDITION} and
\textsc{FALLING\_ACTION} relations are collapsed the same way; and Causal-TimeBank's
CLINK annotations map directly, restricted to the intra-sentence subset (286 of 318
CLINKs, 90\%; the remaining cross-sentence CLINKs are dropped at preprocessing, so
Causal-TimeBank has no genuine inter-sentence positives and consequently no
intra/inter split anywhere in this paper). For MAVEN-ERE, EventStoryLine, and
Causal-TimeBank, \textit{norel} pairs are not separately annotated in the source
data; they are the complement of the positive pairs within $P$.

Document-level inference means the model produces $\ell$ for \emph{all} of $P$ in a
single generation, formatted as a JSON array, in contrast to pair-level inference,
which issues $|P|$ independent queries, one per candidate pair.

We follow \citet{gao2023chatgpt}'s preprocessing and evaluation protocol for
comparability with \citet{hao2026sere} and \citet{cai2025dreci}: micro
precision/recall/F1 over the positive (causal) classes, \textit{norel} excluded from
the denominator, reported on the intra-sentence subset where a corpus distinguishes
intra- from inter-sentence pairs --- the subset SERE itself evaluates, and hence the
fair comparison point. We additionally adopt a binary framing throughout: once a
pair is identified as causal, whether it is scored as \textit{causes} or
\textit{causedby} is collapsed into a single positive class before computing
metrics, since the object of interest in this work is \emph{whether} a causal link
exists, not which of the two mentions is cause versus effect. Directed
(non-collapsed) metrics are logged for every run but are not the primary quantity we
report or optimize for.

% ======================================================================
\section{Method}
\label{sec:method}

\subsection{Overview}
\begin{itemize}
  \item Pipeline (Figure~\ref{fig:overview}): test document $\to$ retrieve $k$
        neighbour training documents $\to$ generate (and cache) synthetic CoT
        for each neighbour $\to$ assemble multi-turn prompt (demonstrations as
        past exchanges, then the test document) $\to$ single document-level LLM
        call $\to$ parse JSON $\to$ optional resampling with majority vote.
  \item Two components carry the contribution: the document-level structured
        prompt (\S\ref{sec:docprompt}) and the synthetic-CoT demonstrations
        (\S\ref{sec:synthcot}); retrieval (\S\ref{sec:retrieval}) and
        aggregation (\S\ref{sec:inference}) are deliberately simple.
\end{itemize}

\begin{figure}[t]
  \centering
  \todo{Figure 1: pipeline overview --- neighbour retrieval, 2-step CoT
  synthesis (gold-guided generation -> decontamination), multi-turn prompt
  assembly, single doc-level call, majority vote.}
  \caption{Overview of \method{}.}
  \label{fig:overview}
\end{figure}

\subsection{Document-Level Structured Prompting}
\label{sec:docprompt}
\begin{itemize}
  \item One call per document; every pair of $P$ requested at once; the model
        sees the full narrative, not a windowed snippet.
  \item In-prompt reasoning protocol (four ordered steps, enforced by the
        system prompt):
        \begin{itemize}
          \item Step 1 --- mention-level pre-filter: affirmatively argued
                exclusions (``$m$ is a location with no agentive role''), never
                blanket dismissal; surviving mentions listed.
          \item Step 2 --- per-pair arguments \emph{for} (connective, mechanism,
                counterfactual asymmetry, discourse structure) and
                \emph{against} (temporal-only, confounder, reporting verb,
                distance).
          \item Step 3 --- qualitative ranking of retained pairs by causal
                plausibility, no labels committed yet.
          \item Step 4 --- final JSON over \emph{every} requested pair;
                everything pre-filtered defaults to \textit{norel}.
        \end{itemize}
  \item Precision-first rule block (targets causal hallucination directly):
        two-signal rule for non-overt links; stricter gate for inter-sentential
        / long-range pairs; default-to-\textit{norel} under uncertainty;
        confounder / proximal-cause check; reporting and temporal connectives
        explicitly ruled non-causal; counterfactual direction test.
  \item In-prompt coherence rules: symmetry
        ($\textit{causes}(i,j) \Leftrightarrow \textit{causedby}(j,i)$) and
        dataset-specific transitivity rules --- checkable at parse time because
        the whole graph is produced jointly \todo{decide whether the coherence
        checker/post-hoc rule pruning appears in main text or appendix}.
  \item Why this leverages better models: the protocol asks for exactly the
        global, long-context reasoning that frontier models are differentially
        good at; pair-level scaffolds cap that ability at the width of one pair.
\end{itemize}

\subsection{Neighbour-Document Retrieval}
\label{sec:retrieval}
\begin{itemize}
  \item Demonstration pool: training split of the same corpus (capped pool
        size to bound CoT-generation cost) \todo{final pool size}.
  \item Selection: TF-IDF cosine over full document text, top-$k$ ($k=3$);
        alternatives evaluated in \S\ref{sec:ablations}: random, mention-set
        Jaccard, sentence-embedding cosine.
  \item Demonstrations are \emph{whole documents with their complete gold
        labelling}, including the \textit{norel} pairs --- every demo carries
        explicit negative evidence, the pair-level analogue of which SERE must
        engineer via balanced retrieval.
\end{itemize}

\subsection{Synthetic Chain-of-Thought Generation}
\label{sec:synthcot}
\begin{itemize}
  \item Motivation: label-only demos teach the output schema, not the
        reasoning; document-level human rationales are unavailable and
        prohibitively expensive to write; gold labels for training documents
        are available --- so rationales can be \emph{synthesized}.
  \item Default 2-step protocol, per demonstration document:
        \begin{itemize}
          \item (i) \emph{gold-guided generation}: the demo document is shown
                under the exact task prompt, plus its gold labels and the
                instruction to write reasoning that ``leads naturally to these
                labels from the text alone'' without revealing foreknowledge;
          \item (ii) \emph{decontamination}: the response is rewritten from
                scratch to purge leakage phrasing (``the labels show\dots'',
                ``as indicated by the expected output\dots''), so each sentence
                reads as analysis, not confirmation.
        \end{itemize}
  \item 3-step variant: (i) blind attempt with no gold access, (ii) gold-based
        correction that fixes wrong conclusions while keeping correct
        reasoning, (iii) decontamination --- preserves the texture and
        difficulty profile of a genuine attempt \todo{report 2- vs.\ 3-step
        comparison in \S\ref{sec:ablations}}.
  \item Injection: each demo becomes a (user task turn, assistant CoT turn)
        exchange preceding the test document --- the model ``remembers having
        done'' $k$ documents correctly, reasoning included; compatible with
        multi-step task prompts (one synthetic response per step).
  \item Generated once per (document, protocol), cached across the run;
        generator model = inference model by default (self-generated
        rationales) \todo{teacher-vs-self generation experiment --- does a
        stronger generator lift a weaker reasoner?}.
  \item Relation to STaR-style rationalization \citep{zelikman2022star}:
        same trick --- condition on the answer to obtain reasoning --- but used
        as in-context demonstrations rather than fine-tuning data, and with an
        explicit decontamination pass, which we show is load-bearing
        (\S\ref{sec:ablations}).
\end{itemize}

\subsection{Inference and Aggregation}
\label{sec:inference}
\begin{itemize}
  \item Temperature 0; strict JSON parsing with bounded retries; missing pairs
        default to \textit{norel}.
  \item Optional resampling: $N$ independent passes with per-pair majority
        vote, ties broken to \textit{norel} (precision-preserving); pair order
        shuffled per pass to break position bias \todo{state whether the main
        table uses resampling and with which $N$}.
  \item Cost: 1 call per document ($N$ with resampling) vs.\ $|P|$ calls for
        pair-level methods --- plus SERE's per-pair retrieval infrastructure
        (ConceptNet + parser) and Dr.ECI's 4 agent calls per pair; quantified
        in \S\ref{sec:cost}.
\end{itemize}

% ======================================================================
\section{Experiments}
\label{sec:experiments}

\subsection{Datasets and Metrics}
\begin{itemize}
  \item EventStoryLine (ESC) \citep{caselli2017event}: 22 topics, 258 docs,
        1{,}770 causal pairs \todo{confirm version used (0.9 vs 1.5) and align
        statement with Gao preprocessing}.
  \item Causal-TimeBank (CTB) \citep{mirza2014annotating}: 183 docs, 318 causal
        pairs --- sparsest positives, hence the dataset where hallucination is
        most punishing; no intra/inter split.
  \item MAVEN-ERE \citep{wang2022mavenere}: large-scale Wikipedia corpus,
        57{,}992 causal pairs \todo{state evaluated subset/sample size}.
  \item Preprocessing follows \citet{gao2023chatgpt}; micro P/R/F1 on the causal
        classes, \textit{norel} excluded.
  \item \todo{decide whether the MECI multilingual evaluation
        \citep{lai2022meci} goes in (appendix?) --- pipeline supports it.}
\end{itemize}

\subsection{Baselines}
\begin{itemize}
  \item Base: task description only; CoT \citep{wei2022chain}: step-by-step
        before answering; Dr.ECI \citep{cai2025dreci}; SERE \citep{hao2026sere}
        --- all pair-level.
  \item Baseline scores on GPT-4o-mini and Gemini-1.5-pro reproduced from
        \citet{hao2026sere} (their Table~2); Dr.ECI PaLM2 row from
        \citet{cai2025dreci}. We do not re-run them.
  \item \todo{consider adding a same-backbone Base/CoT row (our model,
        pair-level) so the doc-level effect is separable from the backbone
        effect --- currently the table conflates them.}
\end{itemize}

\subsection{Implementation Details}
\begin{itemize}
  \item Backbone: \todo{state final backbone(s) --- e.g.\ DeepSeek-V3.2 via
        OpenRouter --- and justify the ``leverage better models'' claim with at
        least one frontier + one mid-tier model}.
  \item $k=3$ demonstrations, TF-IDF selection, 2-step CoT protocol,
        temperature 0 \todo{resampling config for final runs; pool size;
        exact prompt files versioned in the repo}.
  \item Synthetic CoT generation cost: \todo{2 (or 3) extra calls per unique
        demo document, amortized by caching --- give measured token counts}.
\end{itemize}

\subsection{Main Results}
\label{sec:results}

% ----------------------------------------------------------------------
% Single merged results table.
% Baseline rows (PaLM2 / GPT-4o-mini / Gemini-1.5-pro blocks) are
% reproduced from SERE (Hao et al. 2026), Table 2: dataset-level
% (not intra/inter split) scores on ESC, CTB, MAVEN-ERE.
% CTB has no intra/inter split in the original data, so it's directly
% comparable throughout. ESC and MAVEN-ERE baseline scores are
% dataset-level; SCOTECI's ESC/MAVEN-ERE scores are intra-sentence-only
% (marked *) since that's the only setting we have results for so far --
% not yet directly comparable, will reconcile once full-dataset SCOTECI
% numbers land.
% Bold = best score per column across the whole table (all rows).
% ----------------------------------------------------------------------
\begin{table*}
\begin{tabular}{l ccc|ccc|ccc}
\toprule
 & \multicolumn{3}{c|}{ESC} & \multicolumn{3}{c|}{CTB} & \multicolumn{3}{c}{MAVEN-ERE} \\
Method & P & R & F1 & P & R & F1 & P & R & F1 \\
\midrule
\textbf{PaLM2:} & & & & & & & & & \\
\quad Dr.ECI$^\dagger$ & 29.0 & 75.4 & 41.9 & - & - & 13.0 & 22.0 & 80.0 & 34.5 \\
\midrule
\textbf{GPT-4o-mini:} & & & & & & & & & \\
\quad Base          & 28.4 & 82.9 & 42.3 & 5.4  & 80.5 & 10.1 & 23.1 & 91.5 & 36.9 \\
\quad CoT           & 29.4 & 80.6 & 43.1 & 6.3  & 79.6 & 11.6 & 25.4 & 87.9 & 39.4 \\
\quad Dr.ECI        & 31.0 & \textbf{90.2} & 46.1 & 8.2  & 91.2 & 15.1 & 26.0 & 93.0 & 40.7 \\
\quad SERE          & \textbf{48.3} & 51.5 & \textbf{49.9} & 13.8 & 36.3 & 20.0 & 34.5 & 54.6 & \textbf{42.3} \\
\midrule
\textbf{Gemini-1.5-pro:} & & & & & & & & & \\
\quad Base          & 24.0 & 80.8 & 37.0 & 4.7  & \textbf{95.6} & 9.0  & 21.2 & \textbf{95.1} & 34.6 \\
\quad CoT           & 25.1 & 85.4 & 38.7 & 4.7  & 88.5 & 8.9  & 23.3 & 89.2 & 37.0 \\
\quad Dr.ECI        & 27.6 & 82.0 & 41.3 & 7.2  & 88.5 & 13.3 & 23.4 & 91.7 & 37.3 \\
\quad SERE          & 36.5 & 59.3 & 45.2 & 9.9  & 69.9 & 17.4 & 29.9 & 60.2 & 39.9 \\
\midrule
\textbf{Ours:} & & & & & & & & & \\
\quad SCOTECI$^*$   & 40.7 & 70.9 & 51.7 & \textbf{61.1} & 39.6 & \textbf{48.1} & \textbf{40.8} & 59.6 & 48.4 \\
\bottomrule
\end{tabular}
\centering
\small
\caption{Main results. Baseline rows reproduced from \citet{hao2026sere} (their Table 2). Best score per column is in bold across all rows. $\dagger$: results from \citet{cai2025dreci}. CTB has no intra/inter split. For ESC and MAVEN-ERE, baseline scores are dataset-level, while SCOTECI ($^*$) reports intra-sentence-only results for now -- not yet directly comparable on those two columns.}
\label{tab:main-results}
\end{table*}
The results from maven-ere need an update as soon as the full prediction is given
\begin{itemize}
  \item \textbf{CTB (the headline).} F1 51.0 vs.\ 20.0 for the best baseline
        (SERE): +31 F1. Precision 53.2 where no LLM baseline exceeds 13.8 ---
        the degenerate high-recall/low-precision regime is broken, on the
        dataset where positives are rarest and hallucination costliest.
  \item \textbf{MAVEN-ERE.} Best precision in the table (35.5); F1 36.6 trails
        SERE's 42.3 but the comparison is not yet apples-to-apples
        (intra-sentence-only, see caveat) \todo{full-dataset run}.
  \item \textbf{ESC.} 43.2, in the CoT-baseline range, below SERE's 49.9 ---
        same caveat \todo{full-dataset run; if the gap persists, discuss:
        ESC's dense per-document annotation may reward per-pair scrutiny}.
  \item \textbf{Reading across rows.} Base/CoT/Dr.ECI all sit at recall
        $\ge$75 with precision $\le$31; SERE trades recall for precision but
        stays below P=50 everywhere; \method{} is the only method in the
        balanced regime on two of three datasets.
  \item \todo{add significance / resampling variance statement once final runs
        land.}
\end{itemize}

% ======================================================================
\section{Analysis}
\label{sec:analysis}
% Single-table constraint: everything below reported as prose deltas
% ("removing X costs Y F1 on CTB") or as figures, never new tables.

\subsection{Ablations}
\label{sec:ablations}
\begin{itemize}
  \item Remove synthetic CoT (demos = document + gold JSON only): isolates the
        contribution of \emph{reasoning-bearing} demonstrations
        \todo{expected: format retained, precision drops --- run}.
  \item Remove few-shot entirely (structured protocol alone): isolates the
        document-level protocol \todo{run}.
  \item Remove the structured protocol (direct JSON answer, demos kept):
        isolates protocol vs.\ demonstrations \todo{run}.
  \item Selection strategy: random vs.\ TF-IDF vs.\ mention-Jaccard vs.\
        embedding cosine \todo{run; SERE's finding that naive retrieval can
        underperform random is the hypothesis to test at document scale}.
  \item Number of demonstrations $k \in \{1, 3, 5\}$ \todo{run; SERE saw
        degradation past $k=2$ at pair level --- doc-level demos are much
        longer, context dilution may bite sooner}.
  \item 2-step vs.\ 3-step CoT protocol \todo{run}.
  \item Decontamination on/off: does leakage phrasing in demos teach the model
        to assert without reasoning? \todo{run --- this justifies the step}.
  \item Coherence rules on/off (in-prompt and post-hoc) \todo{run}.
\end{itemize}

\subsection{Document-Level vs.\ Pair-Level, Same Backbone}
\label{sec:cost}
\begin{itemize}
  \item Same backbone forced into pair-level inference (one pair per call, same
        rules) vs.\ \method{}: separates the formulation effect from the
        backbone effect \todo{run --- this is the key internal comparison}.
  \item Cost accounting: LLM calls and total tokens per document ---
        pair-level: $|P|$ calls, document re-sent each time; Dr.ECI: up to 4
        agents per pair; \method{}: 1 call (+$2k$ amortized CoT-generation
        calls, cached) \todo{measured numbers + a cost-vs-F1 figure}.
\end{itemize}

\subsection{Scaling with Model Quality}
\begin{itemize}
  \item Claim to substantiate: the document-level formulation \emph{widens} its
        advantage as the backbone improves, because it exposes rather than
        bypasses long-context global reasoning.
  \item Model ladder: frontier closed / large open / mid open / small open
        \todo{pick 4--6 from the OpenRouter matrix; plot F1 (doc-level vs
        pair-level) against model tier --- figure, not table}.
  \item Expected picture: weak models over-generate at document scale
        (precision collapse on dense docs); strong models benefit; the
        crossover is the point of the paper.
\end{itemize}

\subsection{Error Analysis}
\begin{itemize}
  \item False-positive taxonomy: transitive flattening (A$\to$B$\to$C read as
        A$\to$C when the scheme does not credit it), temporal-as-causal,
        reporting verbs \todo{sample and count from traces}.
  \item Dense documents: does precision degrade with the number of gold pairs
        per document? \todo{scatter/figure}.
  \item Direction errors: \textit{causes} vs.\ \textit{causedby} confusions
        as a fraction of true-pair errors \todo{compute from logs}.
  \item Qualitative: one worked example --- synthetic CoT demo, model reasoning,
        prediction vs.\ gold \todo{pick a CTB doc}.
\end{itemize}

% ======================================================================
\section{Conclusion}
\begin{itemize}
  \item Reframed LLM-based ECI from per-pair scaffolding to a single
        document-level call with structured reasoning and coherence rules.
  \item Introduced synthetic, decontaminated chain-of-thought demonstrations:
        rationalization at inference time, no human rationales, no knowledge
        base, no fine-tuning.
  \item Large precision-driven gains where hallucination is costliest (CTB:
        +31 F1); balanced P/R elsewhere.
  \item The method is a bet on model quality: as backbones improve, the same
        prompt gets better --- the scaffold no longer caps the model.
  \item Future work: full-document ESC/MAVEN runs, multilingual MECI,
        per-document strategy routing, teacher-generated CoT for small models.
\end{itemize}

% ======================================================================
\section*{Limitations}
\begin{itemize}
  \item ESC and MAVEN-ERE results are currently intra-sentence only; the
        headline CTB comparison is clean, the other two columns are not yet.
  \item Baseline numbers are reproduced from \citet{hao2026sere}, not re-run
        with our backbone; the same-backbone pair-level control
        (\S\ref{sec:cost}) is the mitigation.
  \item Synthetic CoT costs $2$--$3$ generation calls per demonstration
        document (amortized by caching, but nonzero).
  \item Decontamination is itself LLM-based; residual leakage phrasing cannot
        be ruled out \todo{manual audit of $n$ sampled demos}.
  \item Long documents with large pair sets stress the single-call format
        (output-length limits, attention dilution); per-document routing or
        chunking may be needed at the extreme.
\end{itemize}

% ======================================================================
% \appendix
% \section{Prompts}
% \begin{itemize}
%   \item Full system/user prompts per dataset (versioned in repo).
%   \item CoT-generation prompts: gold hint, correction hint, decontamination.
%   \item Coherence rule sets per dataset.
% \end{itemize}
% \section{Additional Details}
% \begin{itemize}
%   \item Dataset statistics and label mappings.
%   \item Resampling/tie-breaking details; JSON-repair procedure.
%   \item MECI multilingual results (if not in main text).
% \end{itemize}

% acl.sty already sets \bibliographystyle{acl_natbib} internally.
\bibliography{references}

\end{document}
 